\section{Preliminaries}  \label{sec: prelimnares}
In this section we introduce the spherical Hecke algebra and their Pre-canonical bases. 


\subsection{Affine Weyl groups}

From now on we fix a positive integer $n$.
Let $\Phi$ be a root system of type $A_n$. 
We use the following realization of $\Phi$. 
Let $V\subset \mathbb{R}^{n+1}$ be the subspace of vectors with coordinates adding up to zero. 
In this setting we set $\Phi=\{ \epsilon_i -\epsilon_j \mid 1\leq i\neq j \leq n+1\}$. 
The simple roots are $\Delta = \{ \alpha_i\coloneqq \epsilon_i-\epsilon_{i+1} \mid 1\leq i \leq n \}$ and the set of positive roots is given by $\Phi^+=\{   \alpha_{i,j} \coloneqq  \epsilon_i -\epsilon_{j+1} \mid 1\leq i\leq j \leq n  \}$. 
The fundamental weights $\{\varpi_1 , \ldots , \varpi_n\}$ are given by the rule $\pint{\varpi_i}{\alpha_j} = \delta_{i,j}$. 
The weight lattice is $X=\operatorname{span}_{\mathbb{Z}} \{\varpi_1 , \ldots , \varpi_n\}$ and the dominant weights are $X^+=\operatorname{span}_{\mathbb{N}} \{\varpi_1 , \ldots , \varpi_n\}$.



\begin{remark}\rm \label{remark producto interno}
  Let $1\leq i < j \leq n$ and $1\leq k \leq n$. Then, we have
  \begin{equation}
      \pint{\alpha_{i,j}}{\alpha_k} = \begin{cases}
          1, & \mbox{if } k=i \mbox{ or } k=j;\\
          -1, & \mbox{if } k=i-1 \mbox{ or } k=j+1; \\
          0, & \mbox{otherwise.}
       \end{cases}
  \end{equation}
\end{remark}


\begin{definition}\rm 
  We  define a total order $\preceq $ on  $\Phi^{\geq 2}$ as follows. 
  We write $\alpha_{i,j}\preceq \alpha_{k,r}$ if either $j-i < r-k$,  or $j-i = r-k$ and $i\leq k$.   
\end{definition}

Notice that the root $\alpha_{1,n}$ (resp. $\alpha_{1,2}$) is the maximal (resp. minimal) element for this order.
Given $a\preceq b\in \Phi^{\geq 2}$, we define the  closed interval $[a,b]$ to be  the set of elements $x \in \Phi^{\geq 2}$ such that $a\preceq x \preceq b$. 
Open and half-open intervals are defined in a similar fashion.
For $\alpha_{i,j}\in \Phi^{\geq 2}$ we define
\begin{equation}
    \Phi^{\succeq \alpha_{i,j}} =\{ \alpha \in \Phi^{\geq 2} \mid \alpha \succeq \alpha_{i,j}  \} \qquad \mbox{and} \qquad
     \Phi^{\succ \alpha_{i,j}} =\{ \alpha \in \Phi^{\geq 2} \mid \alpha \succ \alpha_{i,j}  \}.
\end{equation}



Henceforth we represent  $\Phi^+$ as an array of boxes in the plane where each box corresponds to a unique positive root.
We draw a ``pyramid'' of boxes with $n$ boxes at the bottom. 
The boxes on the bottom represent the simple roots and they are represented by the corresponding index. 
The other positive roots are obtained as follows. 
The box associated to $\alpha_{i,j}\in \Phi^{\geq 2}$ is the unique box $B$ such that the pyramid with apex $B$ has as corners of its base the boxes associated to the simple roots $\alpha_i$ and $\alpha_j$. 
Given $K\subseteq \Phi^+$ we color with green a box if its corresponding root belongs to $K$. 
In this setting, it is quite easy to compare to roots $\alpha, \beta \in \Phi^{\geq 2}$.
Namely, let $A$ and $B$ be the boxes associated to $\alpha$ and $\beta$, respectively.
Then, $\alpha \succ \beta$ if and only if either $A$ is located above $B$, or $A$ and $B$ are located at the same level and $A$ is on the right of $B$.
For instance, the array of boxes in  \Cref{fig: rootsA} corresponds to $\Phi^{\succ \alpha_{3,6}}=\Phi^{\succeq \alpha_{4,7}}$ for $n=10$. 


\begin{figure}[H]
     \centering
     \begin{subfigure}[b]{0.45\textwidth}
         \centering
        {\scalebox{.4}{\exbonito}}
    \caption{ $\Phi^{\succ \alpha_{3,6}}=\Phi^{\succeq \alpha_{4,7}}$.  }
    \label{fig: rootsA}
     \end{subfigure}
     \hfill
     \begin{subfigure}[b]{0.45\textwidth}
         \centering
        {\scalebox{.4}{\exbonitoB}}
    \caption{$ K \sim_T \Phi^{\succ \alpha_{3,6}} $ for $T=\{4,5\}$.}
    \label{fig: rootsB}
     \end{subfigure}
        \caption{Roots in a pyramid for $n=10$.}
        \label{fig: ejemplos de roots}
\end{figure}




Let $\alpha\in \Phi^+$. We define $s_{\alpha}: V\rightarrow V$ by $s_\alpha(\lambda) = \lambda - \pint{\lambda}{\alpha}\alpha $. 
The (finite) Weyl group $W_f$ of $\Phi$ is the group generated by $\{s_{\alpha}\mid \alpha \in \Phi^+ \}$. 
Let $s_i = s_{\alpha_i}$ and $S_{f}=\{s_i \mid 1\leq i \leq n\}$. 
Then $W_f$ is a Coxeter system with generators $S_f$. 
As usual, we denote by $\ell $ and $\leq$ its length function and Bruhat order, respectively. 
Also, we denote by $w_0$ the longest element of  $W_f$. 
For $\alpha=\alpha_{i,j}\in \Phi^{+}$, the corresponding reflection $s_{\alpha_{i,j}}$ is given by
\begin{equation}\label{eq:reflection}
s_{\alpha_{i,j}} = s_i s_{i+1}\cdots s_{j-1}s_j s_{j-1}\cdots s_{i+1}s_i.
\end{equation}




Besides the natural action of $W_f$ on $V$ we have the \emph{dot} action that is defined as follows. 
Let $\rho $ be the half-sum of the positive roots. 
We have $\rho = \sum_{i=1}^n \varpi_i$. 
The dot action is given by $w\cdot \lambda = w(\lambda +\rho) -\rho$, for $w\in W_f$ and $\lambda \in V$. 

Let $m\in\mathbb{Z}$ and $\alpha\in\Phi^+$. 
Define the (affine) hyperplane
$H_{\alpha,m}:=\{\lambda\in V \mid \pint{\lambda}{\alpha}=m\}$,
and let $s_{\alpha,m}$ denote the reflection along $H_{\alpha,m}$.
In formulas, $s_{\alpha,m}(\lambda)=\lambda-(\pint{\lambda}{\alpha}-m)\,\alpha$.
The affine Weyl group $W_a$ is the group generated by $\{s_{\alpha,m } \mid \alpha\in \Phi^+, m\in \mathbb{Z}\}$. 

The connected components of the complement of the union of all hyperplanes
$H_{\alpha , m}$ are called \emph{alcoves}. 
We denote by $\mathcal{A}$ the set of all alcoves. 
We call 
$$\mathcal{A}_0 = \{ \lambda \in V \mid -1 < \pint{\lambda}{\alpha} < 0  \mbox{ for any } \alpha \in \Phi^+  \}$$ the fundamental alcove. 
The map $W_a \rightarrow \mathcal{A}$ given by $w\mapsto w(\mathcal{A}_0)$ is a bijection. 

The walls of $\mathcal{A}_0$ are supported in the hyperplanes $H_{\alpha , 0}$ for $\alpha \in \Delta$ and $H_{\beta , -1}$ for $\beta = \alpha_{1,n}$. 
We set $s_{0} = s_{\beta , -1}$, the reflection along $H_{\beta , -1}$ and $S_a = S_f \cup \{s_0\}$. 
Then, $W_a$ is a Coxeter group with generators $S_a$. 

\begin{definition}\rm \label{def: theta lambda}
    Let $\lambda \in X^+$. 
    The set $w_0(\mathcal{A}_0)+\lambda$ belongs to $\mathcal{A}$. 
    Then we define $\theta(\lambda) \in W_a$ to be the unique element in $W_a$ such that $\theta(\lambda) (\mathcal{A}_0) =  w_0(\mathcal{A}_0)+\lambda$. 
\end{definition}





\subsection{Spherical Hecke algebras}


Let $\mathcal{H} = \mathcal{H}(W_{a}, S_a)$  be the affine Hecke algebra of $(W_a,S_a)$. 
It is the $\mathbb{Z}[q^{\frac{1}{2}},q^{-\frac{1}{2}}]$-algebra  with generators $\{H_{s_i} \mid s_i \in S_a\}$ and relations 
$H_{s_i}^2 = (q^{-\frac{1}{2}} -q^{\frac{1}{2}})H_{s_i}+1$, $H_{s_i}H_{s_j}=H_{s_j}H_{s_i}$ if $s_i$ and $s_j$ commute, and   $H_{s_i}H_{s_j}H_{s_{i}} =H_{s_j}H_{s_i}H_{s_j}$ if $s_i$ and $s_j$ do not commute. 
The Hecke algebra $\mathcal{H}$ comes equipped with a  standard basis $\{\mathbf{H}_w \mid w\in W_a\}$ and a canonical basis $\{ \underline{\mathbf{H}}_{w} \mid w\in W_a \ \}$ 
(see \cite{Lusztig1979,Soergel1997KazhdanLusztigPA}).
The coefficients in the transition matrix between the canonical and standard bases of $\mathcal{H}$ are the Kazhdan-Lusztig polynomials that we denote by $h_{x,y}(q^{1/2})\in \mathbb{Z}[q^{1/2}]$ for $x, y \in W_a$. 
In formulas, we have
\begin{equation}
    \underline{\mathbf{H}}_y = \sum_{x\in W_a} h_{x,y}(q^{1/2}) \mathbf{H}_x. 
\end{equation}

\begin{definition}\rm
   For $\lambda \in X^+$ we set 
$\underline{\mathbf{H}}_{\lambda} = \underline{\mathbf{H}}_{\theta (\lambda)}$. 
The spherical Hecke algebra $\sph $ is defined as
\begin{equation}
    \sph = \operatorname{span}_{\mathbb{Z}[q^{\frac{1}{2}},q^{-\frac{1}{2}}]} \{ \underline{\mathbf{H}}_{\lambda} \mid \lambda \in X^{+}  \}. 
\end{equation} 
We refer to $ \{ \underline{\mathbf{H}}_{\lambda} \mid \lambda \in X^{+}  \}$ as the canonical basis of $\sph$. 
\end{definition}

\begin{remark}\rm
As its name suggests, $\sph$ is an algebra.
However, this is not a subalgebra of $\mathcal{H}$.
In this paper we use only the linear structure of $\sph$, so we do not recall the product. 
For the ring  structure, see \cite{LPP}.
\end{remark}

The spherical Hecke algebra also has a standard basis $\{ \mathbf{H}_{\lambda} \mid  \lambda \in X^+ \}$, where
\begin{equation}\label{eq: defi standard basis}
     \mathbf{H}_{\lambda} = \sum_{x \in D_{\lambda}} q^{-\frac{1}{2}(\ell(\theta(\lambda)) - \ell(x))} \mathbf{H}_{x},   
\end{equation}
 $D_\lambda= W_f  \,\theta (\lambda) \, W_\lambda$ and 
$W_{\lambda}$ is the maximal parabolic subgroup of $W_{a}$ generated by $S_a\setminus \{s_k\}$ and $0\leq k \leq n$ is the unique integer such that $\lambda \equiv -\varpi_k \mod \mathbb{Z}\Phi$ ($\varpi_0\coloneqq 0$ by convention). 

For the sake of completeness, we recall here the definition of the pre-canonical bases given in \Cref{defi pre-canonical}. 

\begin{definition}\rm
   For $i\geq 2 $ the $i$-th pre-canonical basis, $\mathcal{N}^{i}=\{ \bfN_{\lambda }^i  \mid  \lambda \in X^{+} \}$, is given by 
\begin{equation}  \label{def: pre-canonical}
    \bfN_{\lambda}^{i} = \sum_{I\subset \Phi^{\geq i}} (-q)^{|I|} \widetilde{\underline{\mathbf{H}}}_{\lambda - \Sigma_I}, 
\end{equation}  
As its name suggests, for each $i$, $\mathcal{N}^{i}$ is a basis for $\sph$.
\end{definition}

The main result of this paper is a proof of the positivity statement in \Cref{Conj A}. 
The key idea is that pre-canonical basis elements are too rigid for our purposes: they are defined only for dominant weights, and their construction is constrained to  sets of the form $\Phi^{\ge i}$ for some $i\ge 2$. 
We address these two limitations as follows.

\begin{definition}\rm
      For $A\subset \Phi$ and $\mu \in X$ we define
      \begin{equation}
          \bfM_{\mu}^{A}\coloneqq \sum\limits_{I\subset A} (-q)^{\labs I \rabs}\bfH_{\mu - \Sigma_{I}}
      \end{equation}
    In particular, we define $\bfM_{\lambda}^{\succeq \alpha_{i,j}}= \bfM_{\lambda}^{\Phi^{\succeq \alpha_{i,j}}}$ and  $\bfM_{\lambda}^{\succ \alpha_{i,j}}= \bfM_{\lambda}^{\Phi^{\succ \alpha_{i,j}}}$.
\end{definition}


The following lemma is the key to rewrite $\bfM$-elements.  It is proved in \cite[Proposition 4.3]{LPP}.


\begin{lemma}\rm\label{lem: weyl group over M elements}
Let $A\subset \Phi^+$ and $\lambda\in X$. Then, for all $w\in W_f$ we have 
\begin{equation}\label{eq: lemma act on M-elements}
  \bfM_{\lambda}^{A} = (-1)^{\ell(w)}\,\bfM_{w\cdot \lambda}^{\,w(A)}.
\end{equation}
In particular, if $w(A)=A$, $w\cdot \lambda=\lambda$, and $\ell(w)$ is odd, then $\bfM_{\lambda}^{A}=0$.
Moreover, if $w=s_{\alpha}$ for some $\alpha\in\Phi^{+}$ and $\pint{\lambda+\rho}{\alpha}=0$, then $\bfM_{\lambda}^{A}=0$. 
\end{lemma}

The above lemma has two easy corollaries that shows how we can modify the set indexing a $\bfM$-element, while keeping the same weight. 

\begin{corollary}\label{coro puedo agregar}
    Let $A\subset \Phi^+$ and $\lambda \in X$.  
    Let $ 1\leq r \leq n$ and $\beta \in \Phi^+$.
    Suppose that $s_r(A)=A$,  $\pint{\lambda}{\alpha_r}=0$ and $\pint{\beta}{\alpha_r} =1$ then 
    \begin{equation}
        \bfM_{\lambda}^{A} = \bfM_{\lambda}^{A\cup \{\beta\} }.
    \end{equation}
\end{corollary}

\begin{proof}
By the definition of $\bfM$-elements we have $\bfM_{\lambda}^{A\cup \{\beta\} } = \bfM_{\lambda}^{A} -q \bfM_{\lambda -\beta }^{A}  $. 
Since $\pint{\lambda -\beta +\rho}{\alpha_r}=0$ we conclude by \Cref{lem: weyl group over M elements} that $\bfM_{\lambda -\beta }^{A}=0$, and the result follows. 
\end{proof}

\begin{corollary}\rm\label{lem: sacar una poner otra}
    Let $A\subseteq \Phi^{\geq 2} $ and $1\leq r \leq n$. Suppose that $A\setminus s_r(A)=\{\alpha \}$ and that $ \beta \notin A$. Furthermore, assume that $\langle  \alpha, \alpha_r \rangle = \langle   \beta , \alpha_r  \rangle =1$. If $\lambda \in X$ and $\langle \lambda , \alpha_r \rangle =0$ then 
    \begin{equation}
        \bfM_{\lambda}^A = \bfM_{\lambda}^{(A\setminus \{\alpha\}) \cup \{\beta\}  }. 
    \end{equation}
\end{corollary}

\begin{proof}
   We begin by noticing that $s_r(A\setminus \{ \alpha  \}) = A\setminus \{ \alpha \}$ and $\pint{\lambda -\alpha +\rho}{\alpha_r}=0$. Therefore, \Cref{lem: weyl group over M elements} implies that $\bfM_\lambda^A =\bfM_{\lambda}^{A\setminus \{\alpha\}}$. 
   On the other hand, by applying \Cref{coro puedo agregar} we obtain $\bfM_{\lambda}^{A\setminus \{\alpha \}} = \bfM_{\lambda}^{(A\setminus \{\alpha\}) \cup \{\beta\}  } $.
   By comparing the two expressions we get the desired equality. 
\end{proof}

% Henceforth, we will say that we apply \Cref{lem: sacar una poner otra} at position $r$ to the triple $(A,\alpha, \beta)$ if we are under the hypothesis of the lemma.  





